\chapter{Informed Trading: Glosten--Milgrom and Kyle}
\label{ch:kyle_gm}
\chaptermark{Kyle and Glosten--Milgrom}

Why does the market maker charge a spread even when they have zero
costs, zero inventory, and zero holding period?

Strip the frictions away. No fees, no clearing charges, no
per-message costs; a book-flat starting position; every fill unwound
for free at the same mid-price an instant later. The three-cause
decomposition of \cref{ch:spreads_roll} has silenced two of its
three sources. The naive prediction is a zero spread.

The prediction is wrong. Even in this frictionless world the
rational market maker quotes a strictly positive spread, because of
\emph{adverse selection}: some counterparties know something the
market maker does not, and every fill against such a counterparty
is a small expected loss. The market maker cannot see who is
informed, so they charge everyone a tax wide enough that the gains
on the uninformed cover the losses on the informed. That tax is
the spread.

This chapter formalises two models of the phenomenon.
Glosten--Milgrom (1985) \citeGM{} is sequential: trades arrive one
at a time, the market maker updates beliefs by Bayes, and the
equilibrium spread falls out of zero expected profit on each
trade. Kyle (1985) \citeKyle{} is batched and strategic: a single
informed trader picks trade size optimally knowing their orders
move prices, and a competitive market maker prices the aggregate
order flow by projection. Both models will be referenced
repeatedly in \cref{ch:market_impact,ch:market_making,ch:almgren_chriss}.

\begin{backgroundbox}[Prerequisites]
From earlier in the book:
\begin{itemize}
  \item \Cref{ch:lob}: the limit order book, limit vs market orders,
  makers and takers, the idea that a market maker's job is to
  quote a bid and an ask simultaneously and wait for the flow to
  come to them.
  \item \Cref{ch:participants}: the \textbf{informed vs uninformed}
  taxonomy, the toy adverse-selection calculation of
  \cref{subsec:toy_adverse} with $\alpha = 0.3$, $V \in \{100, 110\}$
  and a zero-spread market maker losing $\$1.5$ per fill, and the
  Prosperity~3 Olivia story as ground truth for the existence of
  informed flow in the wild.
  \item \Cref{ch:price_value}: the formal definitions of $\bidp$,
  $\askp$, $\spreads$, and $\midp$.
  \item \Cref{ch:spreads_roll}: the three-cause decomposition of the
  spread, the Roll estimator, and the explicit forward references
  that pointed to \emph{this} chapter for the formal treatment of
  the adverse-selection component.
\end{itemize}
From outside the book:
\begin{itemize}
  \item \textbf{Bayes' theorem} and \textbf{conditional
  expectation}: given a prior $\Prob(V = v)$ and a likelihood
  $\Prob(\text{observation} \mid V = v)$, the posterior is
  $\Prob(V = v \mid \text{observation}) = \Prob(\text{observation} \mid
  V = v) \Prob(V = v) / \Prob(\text{observation})$, and
  $\E[f(V) \mid \text{observation}] = \sum_v f(v) \Prob(V = v \mid
  \text{observation})$.
  \item Elementary random walks and the fact that conditioning on
  new information updates beliefs about a latent variable.
  \item For Kyle: the \textbf{projection formula} for bivariate
  Gaussians. If $X, Y$ are jointly normal with $\E[X] = \E[Y] = 0$,
  then $\E[X \mid Y] = (\Cov(X, Y) / \Var(Y)) \cdot Y$. This is a
  one-line result; Appendix~A has the derivation if it is rusty.
  \item Basic calculus of one-variable optimisation (taking a
  derivative and setting it to zero) for the insider's
  profit-maximisation step.
\end{itemize}
\end{backgroundbox}

\section{Glosten--Milgrom: the sequential trade model}
\label{sec:gm}

How little machinery is enough to derive a positive spread? The
Glosten--Milgrom model uses four ingredients: (i) a binary latent
asset value, (ii) informed and uninformed traders in a fixed mix,
(iii) a risk-neutral market maker who quotes bid and ask so that
each trade has zero expected profit, and (iv) Bayes' theorem. Out
of those drops the prediction that the spread is strictly positive
whenever $\alpha > 0$, and narrows as the market maker learns from
the trade sequence.

\subsection{Setup}
\label{subsec:gm_setup}

There is a single risky asset with fundamental value $V$, either
high or low:
\[
  V \in \{V_L, V_H\}, \qquad V_L < V_H.
\]
The prior is $\theta_0 \defeq \Prob(V = V_H) \in (0, 1)$. The true
value $V$ is drawn once at the start of the game; only informed
traders learn it.

A fraction $\alpha \in [0, 1]$ of arriving traders are
\textbf{informed}: they know $V$ and trade in the profitable
direction.
\[
  \text{informed trader} \;\Longrightarrow\;
  \begin{cases}
    \text{buys}  & \text{if } V = V_H, \\
    \text{sells} & \text{if } V = V_L.
  \end{cases}
\]
The remaining $1 - \alpha$ are \textbf{uninformed} (also
\textbf{noise} or \textbf{liquidity} traders): they trade for reasons
exogenous to price (rebalancing, cash needs, hedging). Their
buy/sell decisions are independent of $V$ and symmetric:
\[
  \text{uninformed trader} \;\Longrightarrow\;
  \begin{cases}
    \text{buys}  & \text{with probability } 1/2, \\
    \text{sells} & \text{with probability } 1/2.
  \end{cases}
\]
Each trader arrives, looks at the posted quotes, and trades exactly
one unit (GM has nothing to say about order sizing; Kyle below
does).

The market maker is \textbf{risk-neutral} (cares about expected
profit, not variance) and \textbf{competitive} (quotes at the
break-even level: wider and a competitor undercuts, narrower and
they lose money). Competitive risk-neutral pricing means the quoted
ask and bid equal the expected value conditional on the next trade
being a buy or a sell:
\begin{equation}
  \askp_t \;=\; \E[V \mid \text{next trade is a buy}, \mathcal{F}_t],
  \qquad
  \bidp_t \;=\; \E[V \mid \text{next trade is a sell}, \mathcal{F}_t],
  \label{eq:gm_quotes}
\end{equation}
where $\mathcal{F}_t$ is the history of trades up to time $t$.
Everybody observes each realised trade direction, and the market
maker updates $\theta_t \defeq \Prob(V = V_H \mid \mathcal{F}_t)$
by Bayes before re-posting quotes. Only the ordering of trades
matters, not the physical waits between them.

\subsection{Toy example, worked by hand}
\label{subsec:gm_toy}

Before the general formula, do one Bayesian update with concrete
numbers, chosen to match the adverse-selection toy of
\cref{subsec:toy_adverse}. Take
\[
  V_L = 100, \qquad V_H = 110, \qquad \theta_0 = 0.5,
  \qquad \alpha = 0.3.
\]
Five numbered steps.

\subsubsection{Step 1: setup.}
At $t = 0$, $\Prob(V = 110) = 0.5$, so the unconditional expectation is
\[
  \E[V] \;=\; 0.5 \cdot 110 + 0.5 \cdot 100 \;=\; 105.
\]
Quoting both bid and ask at $105$ loses $-\$1.5$ per fill
(\cref{subsec:toy_adverse}). The question is what to quote instead.

\subsubsection{Step 2: compute $\Prob(\text{buy} \mid V)$.}
With probability $\alpha$ the arriving trader is informed (buys iff
$V = V_H$); with probability $1 - \alpha$ uninformed (buys with
probability $1/2$). Therefore
\begin{align*}
  \Prob(\text{buy} \mid V = V_H)
  \;&=\; \underbrace{\alpha \cdot 1}_{\text{informed, sees }V=V_H\text{, buys}}
       \;+\; \underbrace{(1 - \alpha) \cdot \tfrac{1}{2}}_{\text{uninformed, flips a coin}} \\
  \;&=\; 0.3 + 0.7 \cdot 0.5
  \;=\; 0.3 + 0.35
  \;=\; 0.65,
\end{align*}
and by symmetry
\begin{align*}
  \Prob(\text{buy} \mid V = V_L)
  \;&=\; \underbrace{\alpha \cdot 0}_{\text{informed, sees }V=V_L\text{, sells}}
       \;+\; \underbrace{(1 - \alpha) \cdot \tfrac{1}{2}}_{\text{uninformed, flips a coin}} \\
  \;&=\; 0 + 0.35
  \;=\; 0.35.
\end{align*}
A buy is more likely in the high state ($0.65 > 0.35$), which is
what makes a buy informative about $V$. By symmetry,
$\Prob(\text{sell} \mid V_H) = 0.35$ and $\Prob(\text{sell} \mid
V_L) = 0.65$.

\subsubsection{Step 3: compute the posterior $\Prob(V = V_H \mid \text{buy})$.}
The prior is $\theta_0 = 0.5$. The unconditional probability of a
buy under the mixture is
\[
  \Prob(\text{buy})
  \;=\; \theta_0 \cdot \Prob(\text{buy} \mid V_H)
       \;+\; (1 - \theta_0) \cdot \Prob(\text{buy} \mid V_L)
  \;=\; 0.5 \cdot 0.65 + 0.5 \cdot 0.35 \;=\; 0.5.
\]
At a symmetric prior the aggregate buy rate is $1/2$. Bayes gives
\[
  \Prob(V = V_H \mid \text{buy})
  \;=\; \frac{\Prob(\text{buy} \mid V_H) \cdot \Prob(V_H)}{\Prob(\text{buy})}
  \;=\; \frac{0.65 \cdot 0.5}{0.5} \;=\; 0.65.
\]
The buy updates the high-state posterior from $0.5$ to $0.65$; by
symmetry $\Prob(V = V_H \mid \text{sell}) = 0.35$.

\subsubsection{Step 4: compute $\E[V \mid \text{buy}]$ and $\E[V \mid \text{sell}]$.}
\[
  \E[V \mid \text{buy}]
  \;=\; \Prob(V_H \mid \text{buy}) \cdot V_H
       \;+\; \Prob(V_L \mid \text{buy}) \cdot V_L
  \;=\; 0.65 \cdot 110 + 0.35 \cdot 100 \;=\; 106.5.
\]
By \cref{eq:gm_quotes}, this is the market maker's ask:
\[
  \askp_0 \;=\; 106.5.
\]
Symmetrically, the bid is
\[
  \bidp_0 \;=\; \E[V \mid \text{sell}]
  \;=\; 0.35 \cdot 110 + 0.65 \cdot 100 \;=\; 103.5.
\]
The implied spread and mid at the start of the game are
\[
  \spreads_0 \;=\; \askp_0 - \bidp_0 \;=\; 106.5 - 103.5 \;=\; 3.0,
  \qquad
  \midp_0 \;=\; \frac{\askp_0 + \bidp_0}{2} \;=\; 105.0.
\]

\subsubsection{Step 5: interpret.}
The mid is $105.0$, equal to the prior expected value. The spread
of $3.0$ is positive despite zero costs, zero inventory, and zero
holding period, attributable entirely to adverse selection.

Check break-even: the ask equals $\E[V \mid \text{buy}]$, so the
expected P\&L on the next buy is $\askp_0 - \E[V \mid \text{buy}] =
0$, and by symmetry zero on a sell. The $\$1.5$-per-fill loss of
the naive $105$-quote (\cref{subsec:toy_adverse}) has been priced
away.

\begin{intuitionbox}
The ask $106.5$ sits above the prior mean $105$ because \emph{``the
next trade is a buy'' is itself information about $V$}: a buy is
more likely when $V = V_H$, so conditioning on a buy pushes
$\Prob(V_H)$ from $0.5$ to $0.65$ and the expected value from
$105$ to $106.5$. The bid at $103.5$ is below $105$ by symmetry.

The spread is no free lunch. The market maker's expected profit
per trade is zero. Uninformed traders pay $1.5$ away from fair
value on each trade; that loss, summed over the $(1-\alpha)$
uninformed fraction, exactly covers the $\alpha = 0.3$ of informed
trades on which the market maker loses roughly $5$ per trade
(half the width $V_H - V_L$). The spread is a \textbf{tax on the
entire population} used to insure the market maker against the
informed minority.
\end{intuitionbox}

\subsection{The general GM spread formula}
\label{subsec:gm_general}

Let $\theta_t \defeq \Prob(V = V_H \mid \mathcal{F}_t)$ be the
posterior after $t$ trades. The likelihoods of the next trade
direction, in each state, are functions of $\alpha$ alone:
\[
  \Prob(\text{buy} \mid V_H) \;=\; \alpha + \tfrac{1-\alpha}{2}
  \;=\; \tfrac{1 + \alpha}{2},
  \qquad
  \Prob(\text{buy} \mid V_L) \;=\; \tfrac{1-\alpha}{2}.
\]
Symmetrically $\Prob(\text{sell} \mid V_H) = (1 - \alpha)/2$ and
$\Prob(\text{sell} \mid V_L) = (1 + \alpha)/2$. The unconditional
buy probability at time $t$ is
\begin{equation}
  \Prob_t(\text{buy}) \;=\;
     \theta_t \cdot \tfrac{1 + \alpha}{2}
     \;+\; (1 - \theta_t) \cdot \tfrac{1 - \alpha}{2}
     \;=\; \tfrac{1}{2} + \tfrac{\alpha}{2}(2\theta_t - 1).
  \label{eq:gm_pbuy}
\end{equation}
The posterior after a buy at time $t+1$ is
\begin{equation}
  \theta_{t+1}^{\text{buy}}
  \;=\; \frac{\Prob(\text{buy} \mid V_H)\,\theta_t}{\Prob_t(\text{buy})}
  \;=\; \frac{(1 + \alpha)\,\theta_t}{1 + \alpha (2\theta_t - 1)},
  \label{eq:gm_postbuy}
\end{equation}
and after a sell
\begin{equation}
  \theta_{t+1}^{\text{sell}}
  \;=\; \frac{(1 - \alpha)\,\theta_t}{1 - \alpha (2\theta_t - 1)}.
  \label{eq:gm_postsell}
\end{equation}
The ask and bid for the $(t+1)$-st trade are the conditional
expectations of $V$ under these two posteriors:
\begin{align}
  \askp_t \;&=\; \E[V \mid \text{buy}, \mathcal{F}_t]
         \;=\; \theta_{t+1}^{\text{buy}} V_H + (1 - \theta_{t+1}^{\text{buy}}) V_L, \label{eq:gm_ask} \\
  \bidp_t \;&=\; \E[V \mid \text{sell}, \mathcal{F}_t]
         \;=\; \theta_{t+1}^{\text{sell}} V_H + (1 - \theta_{t+1}^{\text{sell}}) V_L. \label{eq:gm_bid}
\end{align}
Substituting \cref{eq:gm_postbuy,eq:gm_postsell} gives the spread
as an explicit function of $(\alpha, \theta_t, V_H, V_L)$:
\begin{equation}
  \spreads_t
  \;=\; \askp_t - \bidp_t
  \;=\; (V_H - V_L) \cdot
     \left[
       \frac{(1+\alpha)\theta_t}{1 + \alpha(2\theta_t - 1)}
       \;-\;
       \frac{(1-\alpha)\theta_t}{1 - \alpha(2\theta_t - 1)}
     \right].
  \label{eq:gm_spread_full}
\end{equation}
The expression has four clean limiting properties.

\paragraph{Limit 1: $\alpha \to 0$.} With no informed traders,
$\Prob(\text{buy} \mid V_H) = \Prob(\text{buy} \mid V_L) = 1/2$,
so a buy (or a sell) is uninformative about $V$, the posteriors
$\theta_{t+1}^{\text{buy}} = \theta_{t+1}^{\text{sell}} = \theta_t$,
and $\askp_t = \bidp_t$. The spread collapses to zero. No informed
flow $\Rightarrow$ no adverse selection $\Rightarrow$ no spread,
exactly as intuition demands.

\paragraph{Limit 2: $\alpha \to 1$.} Every trader is informed, so a
buy perfectly signals $V = V_H$ and $\theta_{t+1}^{\text{buy}} = 1$,
$\theta_{t+1}^{\text{sell}} = 0$. The quotes collapse onto the
revealed value, giving
\[
  \spreads_t \;=\; V_H - V_L,
\]
the maximum. The market maker quotes only at fair value: any
tighter loses, any wider goes unfilled. Dark pools populated
entirely by informed hedge funds exhibit this pathology: spreads
there are large or the venue goes unused.

\paragraph{Limit 3: symmetric prior $\theta_t = 1/2$.} At a
symmetric prior the posteriors simplify. From
\cref{eq:gm_postbuy,eq:gm_postsell} with $\theta_t = 1/2$ and
$2\theta_t - 1 = 0$,
\[
  \theta_{t+1}^{\text{buy}} \;=\; \frac{(1+\alpha) \cdot 1/2}{1 + \alpha \cdot 0}
     \;=\; \frac{1 + \alpha}{2},
  \qquad
  \theta_{t+1}^{\text{sell}} \;=\; \frac{1 - \alpha}{2}.
\]
Plugging into \cref{eq:gm_ask,eq:gm_bid}:
\begin{align*}
  \askp_t \;&=\; \frac{1+\alpha}{2} V_H + \frac{1-\alpha}{2} V_L,
  &
  \bidp_t \;&=\; \frac{1-\alpha}{2} V_H + \frac{1+\alpha}{2} V_L,
\end{align*}
and
\[
  \spreads_t
  \;=\; \askp_t - \bidp_t
  \;=\; \left[\frac{1+\alpha}{2} - \frac{1-\alpha}{2}\right](V_H - V_L)
  \;=\; \alpha (V_H - V_L).
\]
Exactly, not approximately. For $\alpha = 0.3$, $V_H - V_L = 10$
this gives $\spreads_0 = 3$, matching \cref{subsec:gm_toy}. The
symmetric-prior spread $\alpha (V_H - V_L)$ is GM's cleanest
formula: linear in the informed share, linear in the value range.

\paragraph{Limit 4: $\theta_t \to 0$ or $1$.} As the posterior
concentrates, both quotes converge to the certain value and the
spread collapses. Full learning $\Rightarrow$ zero spread
$\Rightarrow$ trading stops.

\begin{keyfactbox}[Glosten--Milgrom spread formula]
In the Glosten--Milgrom (1985) model \citeGM{}, a risk-neutral,
competitive market maker facing a mixture of an $\alpha$-fraction
of informed traders (who know $V \in \{V_L, V_H\}$) and a
$(1-\alpha)$-fraction of symmetric uninformed traders posts
\begin{align*}
  \askp_t &\;=\; \E[V \mid \text{next trade is a buy}, \mathcal{F}_t], \\
  \bidp_t &\;=\; \E[V \mid \text{next trade is a sell}, \mathcal{F}_t].
\end{align*}
The equilibrium spread is strictly positive whenever $\alpha > 0$,
and at a symmetric prior $\theta_0 = 1/2$ with small $\alpha$ it
satisfies
\[
  \spreads_0 \;=\; \alpha (V_H - V_L) \;>\; 0.
\]
The market maker incurs \emph{no} costs, bears \emph{no} inventory
risk, and holds positions for \emph{no} time. The entire spread is
compensation for adverse selection alone.
\end{keyfactbox}

\subsection{How the market maker learns}
\label{subsec:gm_learning}

After each trade the market maker updates $\theta$ via
\cref{eq:gm_postbuy} or \cref{eq:gm_postsell}, re-computes quotes,
and posts them for the next trader. Over many trades $\theta_t$
wanders toward the truth and the quotes track it.

A concrete trajectory with $V_L = 100$, $V_H = 110$, $\alpha = 0.3$,
$\theta_0 = 0.5$, and first five trades $\text{B, B, B, S, S}$:

\begin{center}
\begin{tabular}{cccccc}
\toprule
$t$ & Trade & $\theta_t$ & $\bidp_t$ & $\askp_t$ & $\spreads_t$ \\
\midrule
$0$ & ---    & $0.5000$ & $103.500$ & $106.500$ & $3.000$ \\
$1$ & B      & $0.6500$ & $105.000$ & $107.752$ & $2.752$ \\
$2$ & B      & $0.7752$ & $106.500$ & $108.650$ & $2.150$ \\
$3$ & B      & $0.8650$ & $107.752$ & $109.225$ & $1.472$ \\
$4$ & S      & $0.7752$ & $106.500$ & $108.650$ & $2.150$ \\
$5$ & S      & $0.6500$ & $105.000$ & $107.752$ & $2.752$ \\
\bottomrule
\end{tabular}
\end{center}

For comparison, a longer mostly-buy trajectory BBBBB-S-BBBB (9
buys, one sell at position 6) from the same prior evolves as

\begin{center}
\begin{tabular}{cccccc}
\toprule
$t$ & Trade & $\theta_t$ & $\bidp_t$ & $\askp_t$ & $\spreads_t$ \\
\midrule
$0$  & ---   & $0.5000$ & $103.500$ & $106.500$ & $3.000$ \\
$1$  & B     & $0.6500$ & $105.000$ & $107.752$ & $2.752$ \\
$2$  & B     & $0.7752$ & $106.500$ & $108.650$ & $2.150$ \\
$3$  & B     & $0.8650$ & $107.752$ & $109.225$ & $1.472$ \\
$4$  & B     & $0.9225$ & $108.650$ & $109.567$ & $0.917$ \\
$5$  & B     & $0.9567$ & $109.225$ & $109.762$ & $0.538$ \\
$6$  & S     & $0.9225$ & $108.650$ & $109.567$ & $0.917$ \\
$7$  & B     & $0.9567$ & $109.225$ & $109.762$ & $0.538$ \\
$8$  & B     & $0.9762$ & $109.567$ & $109.870$ & $0.304$ \\
$9$  & B     & $0.9870$ & $109.762$ & $109.930$ & $0.168$ \\
$10$ & B     & $0.9930$ & $109.870$ & $109.962$ & $0.092$ \\
\bottomrule
\end{tabular}
\end{center}

By $t = 10$, $\theta = 0.993$ and the spread has collapsed from
$3$ to below $0.1$, a $30\times$ compression.

Three observations. First, both mid and spread change after every
trade: the mid tracks $\theta_t$ ($\theta_t V_H + (1-\theta_t)V_L$)
and the spread narrows as $\theta_t$ moves away from $1/2$, since
a more confident posterior makes additional trades less
informative.

Second, the posterior trajectory is a \textbf{martingale}:
$\E[\theta_{t+1} \mid \mathcal{F}_t] = \theta_t$. This is a general
fact about Bayesian updating and will be used throughout
\cref{ch:mean_reversion_ou,ch:black_scholes}; it is worth proving
in the GM setting.

\begin{proposition}[The GM posterior is a martingale]
\label{prop:gm_martingale}
In the Glosten--Milgrom model, the sequence $\{\theta_t\}_{t \geq 0}$
of posteriors on $V = V_H$ is a martingale with respect to the
filtration $\mathcal{F}_t$ generated by the observed trades:
\[
  \E[\theta_{t+1} \mid \mathcal{F}_t] \;=\; \theta_t.
\]
\end{proposition}
\begin{proof}
Condition on $\mathcal{F}_t$; $\Prob_t(\text{buy}) +
\Prob_t(\text{sell}) = 1$. Write the expected posterior as a
weighted sum of the two possible updates:
\begin{align*}
  \E[\theta_{t+1} \mid \mathcal{F}_t]
  \;&=\; \Prob_t(\text{buy}) \cdot \theta_{t+1}^{\text{buy}}
       \;+\; \Prob_t(\text{sell}) \cdot \theta_{t+1}^{\text{sell}} \\
  \;&=\; \Prob_t(\text{buy}) \cdot \frac{\Prob(\text{buy} \mid V_H) \theta_t}{\Prob_t(\text{buy})}
       \;+\; \Prob_t(\text{sell}) \cdot \frac{\Prob(\text{sell} \mid V_H) \theta_t}{\Prob_t(\text{sell})} \\
  \;&=\; \Prob(\text{buy} \mid V_H) \theta_t \;+\; \Prob(\text{sell} \mid V_H) \theta_t \\
  \;&=\; \theta_t \cdot \underbrace{\bigl(\Prob(\text{buy} \mid V_H) + \Prob(\text{sell} \mid V_H)\bigr)}_{=\,1} \\
  \;&=\; \theta_t.
\end{align*}
The $\Prob_t(\cdot)$ factors cancel in each update ratio, leaving
a sum that telescopes to $\theta_t \cdot 1 = \theta_t$.
\qedhere
\end{proof}

The quoted mid $\midp_t = \theta_t V_H + (1 - \theta_t) V_L$ is
therefore also a martingale: the market maker has no ex-ante
directional bias. Tradable prices in efficient markets are
martingales for the same reason: any non-martingale forecast is
tradable until the profit disappears.

Third, the trajectory $(0.5, 0.65, 0.7752, 0.8650, 0.7752, 0.65)$
ends where it was four steps earlier. In the binary GM model a
BS-pair (in either order) returns the posterior to its starting
value, a useful sanity check on hand-computed trajectories.
(This is special to the binary case; multi-state models have more
elaborate path dependence.)

\begin{intuitionbox}
\textbf{The spread as the price of ignorance.} A zero-spread GM
market maker loses $\alpha \cdot (V_H - V_L)/2$ per trade to the
informed. The spread $\alpha (V_H - V_L)$ (at $\theta = 1/2$) is
exactly the width so uninformed traders, paying $\alpha(V_H -
V_L)/2$ each way, collectively cover the loss. The spread is the
\textbf{expected loss to the informed, amortised across the entire
population}: a tax, not a profit.

Poker analogy (\cref{ch:participants}): one card-counter at a
table of nine tourists, nobody can tell them apart. The house
charges a rake that covers the card-counter's expected winnings.
Rake is the spread; card-counter is the informed trader; tourists
are noise; house is the GM market maker. Nobody gets cheated,
nobody gets rich, the house stays in business.
\end{intuitionbox}

\subsection{A short historical and modelling aside}
\label{subsec:gm_hist}

\begin{historybox}
Glosten (Columbia) and Milgrom (Stanford) published ``Bid, Ask and
Transaction Prices in a Specialist Market with Heterogeneously
Informed Traders'' in the \emph{JFE} in 1985 \citeGM{}. Against a
decade of debate over positive spreads, in which the received
answer was order-processing and inventory costs, GM (and Kyle, in
the same year's \emph{Econometrica}) showed adverse selection alone
produces a spread. The technical innovation was bringing Bayes' theorem to
quote-setting. Milgrom later won the 2020 Nobel for his auction
work.
\end{historybox}

\section{Kyle (1985): the one-shot strategic insider}
\label{sec:kyle}

What changes when the informed trader chooses their size? GM's
informed traders are non-strategic: they always trade one unit in
the profitable direction. A real informed trader chooses \emph{how
much}, trading less aggressively when their order would move the
price against them. Kyle (1985) captures this as a fixed point: the
insider's optimal size depends on the market maker's price-impact
coefficient, and the market maker's impact depends on how
aggressively the insider trades.

\subsection{Setup}
\label{subsec:kyle_setup}

The model is one-shot: a single round, after which $V$ is revealed
and everybody settles up. (Kyle's original paper has $N$ equispaced
auctions limiting to Brownian motion; the single-period version
captures the intuition and suffices here.)

There is a single risky asset with fundamental value $V$. Unlike
GM, $V$ is continuous: drawn once from a Gaussian prior
\[
  V \;\sim\; \Normal(0, \sigma_V^2).
\]
Three types of agents observe or trade in this market.

First, a single \textbf{strategic insider} privately observes $V$
before the auction and submits a signed market order $\insidx$
(positive = buy), observing $V$ but not $\noiseu$. We restrict to
\emph{linear} strategies
\[
  \insidx(V) \;=\; \beta \, V,
\]
where $\beta > 0$ is to be solved for. (Linearity is an
equilibrium property, not an assumption: the unique Kyle
equilibrium has linear strategies. We state this as a fact.)

Second, \textbf{noise traders} (liquidity traders in Kyle's paper)
submit an aggregate net order
\[
  \noiseu \;\sim\; \Normal(0, \sigma_u^2), \qquad
  \noiseu \perp V.
\]
Only their aggregate $\noiseu$ is observable, not the individual
trades. The noise traders are indispensable: without them every
order would equal $\beta V$, the market maker could invert to
recover $V$ exactly, and the insider's edge would collapse to zero
before any profit was realised, so the insider would not trade
and the market would unravel.

Third, a \textbf{risk-neutral, competitive market maker} observes
total order flow $y \defeq \insidx + \noiseu$ (she cannot separate
insider from noise) and prices each unit of flow to break even in
expectation. Competition drives the markup to the conditional
expectation of $V$ (any higher invites undercutting; any lower
loses money on average):
\[
  p_1(y) \;=\; \E[V \mid y].
\]
Linear pricing rules give
\[
  p_1(y) \;=\; \lamk \, y
\]
for constant $\lamk > 0$, \textbf{Kyle's lambda}. (Again,
linearity is an equilibrium property, taken as given.)

\begin{figure}[ht]
\centering
\begin{tikzpicture}[
  box/.style={draw, rounded corners=3pt, minimum width=2.8cm,
    minimum height=0.9cm, text centered, font=\small, align=center},
  arr/.style={-Stealth, thick},
]
\node[box, fill=yellow!20] (insider) at (0, 0)  {Insider\\knows $v$};
\node[box, fill=gray!15]   (noise)   at (0,-2)  {Noise traders\\flow $u$};
\node[box, fill=orange!20] (agg)     at (4.5,-1) {Aggregate\\flow $y=x+u$};
\node[box, fill=green!15]  (mm)      at (9.6,-1)  {Market maker\\sets $p$};
\draw[arr] (insider) -- node[above, sloped, font=\footnotesize]
    {$x=\beta(v-p)$} (agg);
\draw[arr] (noise)   -- node[below, sloped, font=\footnotesize]
    {$u\sim\mathcal{N}(0,\sigma_u^2)$} (agg);
\draw[arr] (agg) -- node[above, font=\footnotesize, align=center]
    {observes\\$y$ only} (mm);
% feedback route drawn as three axis-aligned segments
\draw[thick] (mm.north) -- (mm.north |- 0,0.9);
\draw[thick] (mm.north |- 0,0.9) --
    node[above, font=\footnotesize]{$p=\lambda y+\mu_0$}
    (insider.north |- 0,0.9);
\draw[arr] (insider.north |- 0,0.9) -- (insider.north);
\end{tikzpicture}
\caption{Kyle (1985) information flow. The insider submits order size $x$
proportional to the gap between fundamental value $v$ and the current price.
Noise-trader flow $u$ provides cover. The market maker, observing only
aggregate flow $y$, sets price via the linear rule $p = \lambda y + \mu_0$.
The feedback arrow shows the price updating as the market maker extracts the signal.}
\label{fig:kyle_flow}
\end{figure}

$\lamk$ and $\beta$ must be mutually consistent:
\begin{itemize}
  \item Given $\lamk$, $\beta$ maximises insider profit.
  \item Given $\beta$, $\lamk$ is the coefficient of the linear
  projection of $V$ onto $y$.
\end{itemize}
We solve the fixed point in two steps: insider best response to
arbitrary $\lamk$, then Bayes-consistency of $\lamk$ given $\beta$.

\subsection{Worked toy example, before the derivation}
\label{subsec:kyle_toy}

Before the general derivation, work through one realisation. Take
\[
  \sigma_V \;=\; 1, \qquad \sigma_u \;=\; 1.
\]
\Cref{subsec:kyle_deriv} will show the unique linear equilibrium has
\[
  \lamk^\ast \;=\; \frac{\sigma_V}{2\sigma_u} \;=\; \frac{1}{2},
  \qquad
  \beta^\ast \;=\; \frac{\sigma_u}{\sigma_V} \;=\; 1.
\]
Accept these pending derivation: $\insidx = V$ and $p_1 = 0.5 y$.
Suppose $V = 0.5$ and $\noiseu = 0.2$. Then
\begin{align*}
  \insidx \;&=\; \beta^\ast V \;=\; 1 \cdot 0.5 \;=\; 0.5, \\
  y \;&=\; \insidx + \noiseu \;=\; 0.5 + 0.2 \;=\; 0.7, \\
  p_1 \;&=\; \lamk^\ast \cdot y \;=\; 0.5 \cdot 0.7 \;=\; 0.35.
\end{align*}
The market maker clears at $0.35$ (above the prior mean $0$, since
positive flow signals high $V$). Insider profit is
\[
  (V - p_1) \cdot \insidx \;=\; (0.5 - 0.35) \cdot 0.5 \;=\; 0.075.
\]
A few more realisations:

\begin{center}
\begin{tabular}{rrrrrr}
\toprule
$V$ & $\noiseu$ & $\insidx = V$ & $y = \insidx + \noiseu$ & $p_1 = 0.5 y$ & $(V - p_1)\insidx$ \\
\midrule
$+0.5$ & $+0.2$ & $+0.50$ & $+0.70$ & $+0.350$ & $+0.0750$ \\
$+0.5$ & $-0.3$ & $+0.50$ & $+0.20$ & $+0.100$ & $+0.2000$ \\
$-1.0$ & $+0.4$ & $-1.00$ & $-0.60$ & $-0.300$ & $+0.7000$ \\
$+1.5$ & $-0.6$ & $+1.50$ & $+0.90$ & $+0.450$ & $+1.5750$ \\
$\phantom{+}0.0$ & $+1.0$ & $\phantom{+}0.00$ & $+1.00$ & $+0.500$ & $\phantom{+}0.0000$ \\
\bottomrule
\end{tabular}
\end{center}

Every realised profit here is non-negative because the insider
trades with $V$. The derivation below (\cref{subsec:kyle_deriv})
gives expected profit
\[
  \E[\pi] \;=\; \frac{\sigma_V \sigma_u}{2}
  \;=\; 0.5
\]
per round, verified numerically at the end of the chapter.

\subsection{Derivation of the linear equilibrium}
\label{subsec:kyle_deriv}

Three steps: (i) insider expected profit as a function of $\beta$
and $\lamk$, (ii) maximise over $\beta$, (iii) impose
Bayes-consistency of $\lamk$ and solve for
$(\beta^\ast, \lamk^\ast)$.

\paragraph{Step 1: insider expected profit.}
Given $p_1 = \lamk y$, realised profit on a trade of size $\insidx$ is
\[
  \pi \;=\; (V - p_1) \cdot \insidx
  \;=\; (V - \lamk(\insidx + \noiseu)) \cdot \insidx
  \;=\; V \insidx - \lamk \insidx^2 - \lamk \noiseu \insidx.
\]
The insider chooses $\insidx = \beta V$. Conditional on $V$, using
$\E[\noiseu] = 0$:
\[
  \E[\pi \mid V]
  \;=\; V \cdot \beta V \;-\; \lamk (\beta V)^2 \;-\; \lamk \cdot \underbrace{\E[\noiseu]}_{=\,0} \cdot \beta V
  \;=\; \beta V^2 - \lamk \beta^2 V^2.
\]
Now take expectation over $V$, using $\E[V^2] = \sigma_V^2$:
\begin{equation}
  \E[\pi] \;=\; \beta \sigma_V^2 \;-\; \lamk \beta^2 \sigma_V^2
  \;=\; (\beta - \lamk \beta^2)\,\sigma_V^2.
  \label{eq:kyle_profit}
\end{equation}
Quadratic in $\beta$ with positive linear and negative quadratic
terms: unique maximum at finite positive $\beta$.

\paragraph{Step 2: insider best response.}
Differentiate \cref{eq:kyle_profit} in $\beta$ and set to zero:
\[
  \frac{\partial \E[\pi]}{\partial \beta}
  \;=\; \sigma_V^2 - 2 \lamk \beta \sigma_V^2 \;=\; 0.
\]
Dividing by $\sigma_V^2$ (strictly positive by assumption),
\[
  1 - 2 \lamk \beta \;=\; 0
  \quad\Longrightarrow\quad
  \boxed{\;\beta \;=\; \frac{1}{2\lamk}.\;}
\]
The SOC $\partial^2 \E[\pi]/\partial \beta^2 = -2 \lamk \sigma_V^2
< 0$ confirms a maximum. Double the impact, halve the order size.

\paragraph{Step 3: market maker's Bayes-consistent $\lamk$.}
The market maker observes $y = \beta V + \noiseu$ and posts $p_1 =
\E[V \mid y]$. For bivariate zero-mean Gaussians,
\[
  \E[V \mid y]
  \;=\; \frac{\Cov(V, y)}{\Var(y)} \cdot y.
\]
Compute (using $\noiseu \perp V$):
\[
  \Cov(V, y) \;=\; \beta \sigma_V^2,
  \qquad
  \Var(y) \;=\; \beta^2 \sigma_V^2 + \sigma_u^2.
\]
Therefore
\begin{equation}
  \lamk \;=\; \frac{\Cov(V, y)}{\Var(y)}
      \;=\; \frac{\beta \sigma_V^2}{\beta^2 \sigma_V^2 + \sigma_u^2}.
  \label{eq:kyle_mm}
\end{equation}

\paragraph{Step 4: solve the fixed point.}
Two equations, two unknowns:
\begin{align}
  \beta \;&=\; \frac{1}{2\lamk}, \label{eq:kyle_foc} \\
  \lamk \;&=\; \frac{\beta \sigma_V^2}{\beta^2 \sigma_V^2 + \sigma_u^2}. \label{eq:kyle_mmrepeat}
\end{align}
Substitute $\beta = 1/(2\lamk)$ into \cref{eq:kyle_mmrepeat}.
Numerator: $\sigma_V^2/(2\lamk)$. Denominator:
\[
  \beta^2 \sigma_V^2 + \sigma_u^2
  \;=\; \frac{\sigma_V^2}{4\lamk^2} + \sigma_u^2
  \;=\; \frac{\sigma_V^2 + 4 \lamk^2 \sigma_u^2}{4\lamk^2}.
\]
So the RHS is
\[
  \frac{\sigma_V^2/(2\lamk)}{(\sigma_V^2 + 4\lamk^2\sigma_u^2)/(4\lamk^2)}
  \;=\; \frac{2\lamk \sigma_V^2}{\sigma_V^2 + 4\lamk^2 \sigma_u^2}.
\]
Setting this equal to $\lamk$ and cancelling $\lamk > 0$:
\[
  1 \;=\; \frac{2 \sigma_V^2}{\sigma_V^2 + 4\lamk^2 \sigma_u^2}
  \quad\Longrightarrow\quad
  \sigma_V^2 + 4\lamk^2 \sigma_u^2 \;=\; 2 \sigma_V^2
  \quad\Longrightarrow\quad
  4\lamk^2 \sigma_u^2 \;=\; \sigma_V^2.
\]
Therefore
\[
  \lamk^2 \;=\; \frac{\sigma_V^2}{4 \sigma_u^2}
  \quad\Longrightarrow\quad
  \boxed{\;\lamk^\ast \;=\; \frac{\sigma_V}{2\sigma_u}\;}
\]
(positive root). Then
\[
  \boxed{\;\beta^\ast \;=\; \frac{1}{2 \lamk^\ast}
  \;=\; \frac{1}{2 \cdot \sigma_V/(2\sigma_u)}
  \;=\; \frac{\sigma_u}{\sigma_V}.\;}
\]
The equilibrium: insider trades $\insidx^\ast(V) =
(\sigma_u/\sigma_V)V$, market maker prices $p_1(y) = (\sigma_V /
2\sigma_u)y$.

\paragraph{Step 5: expected insider profit.}
Plug into \cref{eq:kyle_profit}:
\begin{align*}
  \E[\pi^\ast]
  \;&=\; (\beta^\ast - \lamk^\ast (\beta^\ast)^2)\,\sigma_V^2 \\
  \;&=\; \left(\frac{\sigma_u}{\sigma_V}
           \;-\; \frac{\sigma_V}{2\sigma_u}\cdot \frac{\sigma_u^2}{\sigma_V^2}\right)\sigma_V^2 \\
  \;&=\; \left(\frac{\sigma_u}{\sigma_V} - \frac{\sigma_u}{2\sigma_V}\right)\sigma_V^2 \\
  \;&=\; \frac{\sigma_u}{2\sigma_V}\,\sigma_V^2
  \;=\; \frac{\sigma_V \sigma_u}{2}.
\end{align*}
Expected profit per round is $\sigma_V \sigma_u / 2$, matching
\cref{subsec:kyle_toy}.

\paragraph{Step 6: how much does the market maker learn?}
Residual uncertainty about $V$ after observing $y$:
\[
  \Var(V \mid y)
  \;=\; \sigma_V^2 - \frac{\Cov(V, y)^2}{\Var(y)}
  \;=\; \sigma_V^2 - \frac{(\beta^\ast \sigma_V^2)^2}{(\beta^\ast)^2 \sigma_V^2 + \sigma_u^2}.
\]
With $\beta^\ast = \sigma_u/\sigma_V$: numerator $\sigma_u^2
\sigma_V^2$, denominator $2\sigma_u^2$, so
\[
  \Var(V \mid y)
  \;=\; \sigma_V^2 - \frac{\sigma_u^2 \sigma_V^2}{2\sigma_u^2}
  \;=\; \sigma_V^2 - \frac{\sigma_V^2}{2}
  \;=\; \frac{\sigma_V^2}{2}.
\]
A single Kyle round erases exactly half the prior variance. The
market maker learns half the truth per round; the insider's
private information dissipates at a corresponding rate.

\subsection{Applying Kyle to the toy example}
\label{subsec:kyle_apply_toy}

With $\sigma_V = \sigma_u = 1$ from \cref{subsec:kyle_toy}:
\[
  \lamk^\ast \;=\; \frac{1}{2 \cdot 1} \;=\; \tfrac{1}{2},
  \qquad
  \beta^\ast \;=\; \frac{1}{1} \;=\; 1,
\]
with expected profit $1/2$. For $V = 0.5$, $\noiseu = 0.2$:
\begin{align*}
  \insidx \;&=\; 1 \cdot 0.5 \;=\; 0.5, \\
  y \;&=\; 0.5 + 0.2 \;=\; 0.7, \\
  p_1 \;&=\; 0.5 \cdot 0.7 \;=\; 0.35, \\
  (V - p_1) \insidx \;&=\; (0.5 - 0.35) \cdot 0.5 \;=\; 0.075.
\end{align*}
Matching \cref{subsec:kyle_toy}. A Monte Carlo over $2 \times 10^5$
draws from $V, \noiseu \sim \Normal(0, 1)$ gives sample mean profit
$0.501$, within $0.2\%$ of the analytical $0.5$.

\begin{keyfactbox}[Kyle's linear equilibrium]
In the Kyle (1985) model \citeKyle{} with prior $V \sim
\Normal(0, \sigma_V^2)$, noise-trade aggregate $\noiseu \sim
\Normal(0, \sigma_u^2)$ independent of $V$, a single strategic
insider submitting $\insidx = \beta V$, and a risk-neutral
competitive market maker observing total order flow $y =
\insidx + \noiseu$ and quoting $p_1 = \lamk y$, the unique linear
equilibrium is
\[
  \boxed{\;\lamk^\ast \;=\; \frac{\sigma_V}{2\sigma_u},
  \qquad
  \beta^\ast \;=\; \frac{\sigma_u}{\sigma_V}.\;}
\]
The insider's expected profit per round is $\sigma_V \sigma_u / 2$,
and the post-trade residual variance $\Var(V \mid y) = \sigma_V^2
/ 2$, exactly half the prior. Kyle's $\lamk^\ast$ has units of
price per unit order flow and is the natural measure of
\textbf{price impact} in the model. It is what every market-impact
study of the last forty years tries to estimate
\citeHasbrouck{}, \citeBouchaud{}.
\end{keyfactbox}

\begin{intuitionbox}
\textbf{The insider hides inside the noise.} $\beta^\ast =
\sigma_u / \sigma_V$ says insider size is proportional to noise
volatility and inversely proportional to signal volatility. Both
are right on reflection.

Loud noise ($\sigma_u$ large): the insider's $\beta V$ gets
absorbed into the dispersion of $y$ and barely moves $p_1$, so
trade large. Quiet noise ($\sigma_u \to 0$): $y \approx \beta V$
and the market maker inverts to $V$, collapsing the edge into
$p_1$; trade small.

A stronger signal ($\sigma_V$ large) means each unit of $V$ is a
bigger secret: trading one extra unit gives away more of $V$, so
trade less. The equilibrium $\beta^\ast = \sigma_u/\sigma_V$ is
the compromise between profiting on a correct $V$ and avoiding
self-impact.

The lesson: \textbf{never trade large in a quiet market}. In
Prosperity, ``Olivia''-style informed flow in a sleepy product is
the small-$\sigma_u$ Kyle case. A team following her aggressively
advertises the signal; following quietly captures the edge.
\end{intuitionbox}

\subsection{A sketch of the multi-period Kyle model}
\label{subsec:kyle_multi}

The one-shot model omits how private information is gradually
incorporated into the price across rounds. Kyle's 1985 paper
\citeKyle{} works out the $N$-round version and takes the
continuous-time limit; a qualitative sketch suffices here.

In $N$ sequential auctions, the insider submits $\insidx_t$ linear
in the residual information $V - p_{t-1}$, and the market maker
posts a price with coefficient $\lamk_t$. Trade too aggressively
early and the signal reveals; too timidly and the rounds run out.

The equilibrium has deterministic sequences $\{\beta_t\},
\{\lamk_t\}$ (see \citeCartea{} for a modern treatment). Headlines:
\begin{enumerate}
  \item Information is revealed gradually: $\Var(V \mid
  \mathcal{F}_t)$ decreases geometrically to zero, linearly in the
  continuous-time limit.
  \item $\lamk_t$ is \emph{constant in time} to leading order: a
  marginal share moves the price the same in round 1 as in round
  $N$. The insider trades more in later rounds to compensate.
  \item The aggregate flow $y_t$ converges to a Brownian motion
  in the continuous-time limit, and the price walks smoothly to
  the true $V$.
\end{enumerate}
The constant-$\lamk$ result justifies using a single Kyle-style
scalar in optimal-execution models (\cref{ch:almgren_chriss}).

\begin{historybox}
Albert ``Pete'' Kyle wrote his 1985 \emph{Econometrica} paper
``Continuous Auctions and Insider Trading'' \citeKyle{} at
Princeton, appearing alongside Glosten--Milgrom. Kyle's headline
result is the continuous-time limit of his $N$-auction game,
where the insider trades continuously at a rate proportional to
$(V - p_t)\,dt$. The single-period version is the classroom
workhorse. The coefficient $\lamk$ is universally known as
``Kyle's lambda,'' and market-impact software across the
industry contains a function for estimating it.
\end{historybox}

\section{Connecting Glosten--Milgrom and Kyle}
\label{sec:gm_vs_kyle}

GM and Kyle approach the same phenomenon from opposite sides:

\begin{center}
\begin{tabular}{>{\raggedright\arraybackslash}p{3.3cm}
                >{\raggedright\arraybackslash}p{5.4cm}
                >{\raggedright\arraybackslash}p{5.4cm}}
\toprule
Feature & Glosten--Milgrom & Kyle (1985) \\
\midrule
Latent asset value           & Binary $V \in \{V_L, V_H\}$ & Gaussian $V \sim \Normal(0, \sigma_V^2)$ \\
Trading protocol             & Sequential, one-at-a-time & Batched, single auction \\
Informed trader              & Myopic, non-strategic & Strategic, size-optimising \\
Noise component              & Symmetric Bernoulli coin-flip & Gaussian aggregate $\noiseu$ \\
Market maker action          & Quote bid and ask each period & Clear a single price at $y$ \\
Price impact summary         & $\spreads$ in price units & $\lamk$ in price per unit flow \\
Learning per trade           & Bayesian update of $\theta_t$ & Single-shot Gaussian projection \\
Limit shapes                 & $\alpha \to 0 \Rightarrow$ spread $\to 0$ & $\sigma_u \to 0 \Rightarrow \lamk \to \infty$ \\
Primary output               & Bid/ask quotes as functions of history & Price impact coefficient \\
\bottomrule
\end{tabular}
\end{center}

Both start with an informed/uninformed mixture and a
zero-expected-profit market maker, and both derive a strictly
positive price-impact function whose width depends on the
informed share and on signal/noise volatilities. They differ in
protocol: GM one-at-a-time, Kyle simultaneous batch. Real markets
live between: exchanges match trades tick-by-tick (GM-flavoured)
but institutional flow arrives in blocks (Kyle-flavoured). A
full treatment uses GM for tick-level updates and Kyle for
block-sizing.

Generalised, both predict a \textbf{square-root law}: $\Delta p
\propto \sqrt{Q}$ for a large trade of size $Q$. Kyle's
single-round result is linear ($\Delta p = \lamk Q$), but
multi-period multi-agent versions collapse to the empirical
square-root shape; \cref{ch:market_impact} covers this.

\begin{pitfallbox}[Kyle's $\lamk$ is not the GM spread]
A common first-reading error is to identify $\lamk$ with the
spread. They have different units.

GM's spread $\spreads$ has units of \emph{price}: the bid--ask gap,
independent of trade size (GM assumes one unit per trade).

Kyle's $\lamk$ has units of \emph{price per unit order flow}: a
sensitivity. There is no bid/ask in pure Kyle; every trade
clears at $p_1(y) = \lamk y$.

Dimensionally, to compare, multiply $\lamk$ by a typical order
size $Q$:
\[
  \text{spread} \;\sim\; 2 \lamk \cdot Q,
\]
a mnemonic, not a theorem. Do not write ``$\lamk$ is the
spread''; it is a category error comparing a sensitivity to a
price.
\end{pitfallbox}

\begin{prosperitybox}
\textbf{Olivia maps to GM.} The Prosperity~3 Olivia story
(\cref{sec:informed_vs_uninformed}) is a GM-world phenomenon.
Olivia was a bot on \texttt{SQUID\_INK} whose buys/sells
preceded fair-value moves in her direction. Top-ten teams
identified her via the trader-ID field in historical logs and
ran the GM market maker's optimal strategy \emph{in reverse}:
instead of widening the spread, they copied her and became
informed themselves.

Calibrate: $V \in [1990, 2010]$ (so $V_H - V_L \approx 20$) and
$\alpha \approx 0.05$ (one bot among twenty). At symmetric prior,
\[
  \spreads_{\mathrm{GM}} \;\approx\; \alpha (V_H - V_L)
  \;\approx\; 1.0 \text{ seashells},
\]
near the observed Round~1 \texttt{SQUID\_INK} spreads. Teams
copying Olivia recouped this tax; teams ignoring her paid it.

\textbf{A large basket order maps to Kyle.} Trading a $100$-unit
basket in Round 2 with Kyle impact $\lamk \approx 0.01$ shifts
the executed price by
\[
  \Delta p \;\approx\; \lamk \cdot 100 \;=\; 1 \text{ seashell}
\]
against the mid. The expected edge must exceed $1$ seashell per
unit or Kyle impact eats it. This is why top teams slice large
orders into small blocks.
\end{prosperitybox}

\begin{gsbox}
On a Goldman Sachs sell-side desk, Kyle's $\lamk$ is an everyday
object.

\textbf{Execution-algorithm calibration.} Execution algos (VWAP,
POV, IS; derived from first principles in
\cref{ch:almgren_chriss}) rest on a linear-impact model $\Delta p
= \eta \cdot (\text{trade rate})$, where $\eta$ generalises
$\lamk$. An execution Strat re-estimates it daily from
parent-order fills and monitors residuals for regime changes.

\textbf{Client-quote pricing on large trades.} A block-desk
client asks ``sell $5$M XYZ, what's my price?'' The pricing tool
runs a Kyle-flavoured impact model, estimates the expected shift
over the unwind horizon, and adds it as a markup/markdown to the
mid, exactly the $\lamk \cdot Q$ arithmetic. Informed clients
(hedge funds with history) get $\lamk$ adjusted upward;
rebalancing pension funds get it adjusted downward.

\textbf{Delta-hedging equity-derivative positions.} Scheduling an
$N$-share delta hedge balances speed against self-impact. With
$N = 200{,}000$ and $\lamk \approx 0.0002$, the naive one-shot
hedge moves the price by $\approx 40$ bp, enough to turn a
structured-product trade from a winner into a loser. This is
why \cref{ch:almgren_chriss} is a central Strats chapter.

\textbf{Research output.} Junior Strats deliver monthly
price-impact refreshes: last six weeks of parent-order logs, a
Kyle-style linear (or square-root) regression on post-trade
reversion, sensitivities reported to senior traders and risk.
The form is textbook Kyle; the values are data-driven.
\end{gsbox}

\section{Key facts to memorise}
\label{sec:ch5_keyfacts}

\begin{itemize}
  \item \textbf{Glosten--Milgrom setup.} Binary $V \in \{V_L, V_H\}$,
  prior $\theta_0 = \Prob(V = V_H)$, informed fraction $\alpha$,
  uninformed traders flip symmetric coins, risk-neutral competitive
  market maker quotes bid and ask as conditional expectations.
  \item \textbf{Glosten--Milgrom spread formula.} $\askp_t = \E[V \mid
  \text{buy}, \mathcal{F}_t]$, $\bidp_t = \E[V \mid \text{sell},
  \mathcal{F}_t]$. At symmetric prior with small $\alpha$,
  $\spreads_0 = \alpha (V_H - V_L) > 0$, and the spread is strictly
  positive for every $\alpha > 0$.
  \item \textbf{Glosten--Milgrom toy numbers.} With $V_L = 100$,
  $V_H = 110$, $\theta_0 = 0.5$, $\alpha = 0.3$: $\bidp_0 =
  103.5$, $\askp_0 = 106.5$, $\spreads_0 = 3.0$. Zero costs, zero
  inventory, $\spreads > 0$.
  \item \textbf{Glosten--Milgrom Bayesian update.} After a buy,
  $\theta_{t+1} = (1+\alpha)\theta_t / (1 + \alpha(2\theta_t - 1))$;
  after a sell, $\theta_{t+1} = (1-\alpha)\theta_t /
  (1 - \alpha(2\theta_t - 1))$. The posterior is a martingale.
  \item \textbf{Spread as tax.} The Glosten--Milgrom spread is the
  expected loss to informed flow, amortised across all traders.
  Uninformed traders cover, collectively, the market maker's
  expected losses to the informed.
  \item \textbf{Kyle setup.} $V \sim \Normal(0, \sigma_V^2)$,
  $\noiseu \sim \Normal(0, \sigma_u^2)$ independent, strategic
  insider submits $\insidx = \beta V$, risk-neutral competitive
  market maker observes $y = \insidx + \noiseu$ and quotes $p_1 =
  \lamk y$.
  \item \textbf{Kyle linear equilibrium.} $\lamk^\ast = \sigma_V /
  (2 \sigma_u)$ and $\beta^\ast = \sigma_u / \sigma_V$.
  \item \textbf{Kyle at $\sigma_V = \sigma_u = 1$.} $\lamk^\ast =
  1/2$, $\beta^\ast = 1$. Realisation $V = 0.5$, $\noiseu = 0.2$:
  insider trades $0.5$, market maker sees $y = 0.7$, clears at
  $p_1 = 0.35$, insider profit $= 0.075$.
  \item \textbf{Expected insider profit.} $\E[\pi^\ast] =
  \sigma_V \sigma_u / 2$ per round, proportional to the product
  of the two volatilities. (Monte Carlo check: $0.501$ over
  $2\times 10^5$ draws, within $0.2\%$.)
  \item \textbf{Residual variance.} $\Var(V \mid y)^\ast =
  \sigma_V^2 / 2$. A single Kyle round erases exactly half the
  prior uncertainty about $V$.
  \item \textbf{Kyle intuition.} $\beta^\ast = \sigma_u/\sigma_V$
  means the insider hides inside the noise. Larger noise
  $\Rightarrow$ larger insider trades; larger signal $\Rightarrow$
  smaller insider trades. Never trade large in a quiet market.
  \item \textbf{Category error.} Kyle's $\lamk$ has units price /
  order flow; the Glosten--Milgrom spread has units of price.
  They are not the same quantity. Related by the dimensional
  heuristic $\spreads \sim 2 \lamk \cdot Q$ for typical order
  size $Q$, but not equal.
  \item \textbf{Cross-reference map.} Empirical generalisation
  and the square-root impact law: \cref{ch:market_impact}.
  Inventory-aware market-making on top of adverse selection:
  \cref{ch:market_making}. Optimal execution with Kyle-style
  impact: \cref{ch:almgren_chriss}.
\end{itemize}

\bigskip

GM and Kyle together take ``spreads exist because of adverse
selection'' from rhetoric to arithmetic, deriving a positive
price of liquidity purely from informational asymmetry. The next
chapter turns to the empirical footprint: the square-root
market-impact law, the most robust regularity in microstructure
and cleanly generated by Kyle's $\lamk$.
